% metrics.tex -- Methodology: metrics (draft 4 Oct 2026). The definitions are the ones used by the
% analysis tools (multi_drone_control/tools/method_metrics.py, push_metrics.py).
% 9 Oct (structure audit): the unused "carried fraction" metric cut; "every Gazebo case is flown at
% least twice" replaced, since several cells are single runs (marked (1)).

\subsubsection{Metrics}
\label{sec:metrics}

All metrics are computed from the logged motion-capture poses of the payload and the UAVs, with the
events (TAKEOFF, the push, DETACH, LAND) from the run's event log.
\begin{description}
    \item[Height error] $e_z = z_L - z_L^{\mathrm{ref}}$, positive with the payload high; in hover, averaged
    over the 15\,s before LAND. \textbf{Heave:} its peak-to-peak over the hover window.
    \item[Horizontal error] $\|\mathbf{e}_{xy}\|$, the horizontal distance between the payload's centre and
    its reference; on a trajectory, its mean over the first lap or the whole figure, and its maximum.
    \item[Payload tilt] The angle between the ring's axis and the vertical: its mean over the hover and its
    peak after an event.
    \item[Disturbance response] For a constant force applied for a fixed time, the displacement along the
    force relative to its mean over the 5\,s before: steady (mean over the last 5\,s of the force), peak
    (while the force acts), overshoot (the largest excursion the other way after it ends) and time to
    recover (after which it stays within 1\,cm). \textbf{Stiffness:} the force over the steady displacement
    with no integral action.
% SETTLED 10 Oct (was a CHECK): aligned with what Table tab:detach and Figure fig:detach measure
% (tools/thesis_sim_tables.py detach_stats, multi_drone_control): peak in the 5 s after DETACH; back in
% band = the tilt stays under the hover's own maximum over the 5 s before plus 0.5 deg; height change =
% largest |z - mean of the 5 s before| in the 12 s after. Was: "the time until the tilt stays under
% 2 deg, and the payload's height loss" (the pre-twin layout study's wording).
    \item[Detach] The peak payload tilt in the 5\,s after DETACH; the time until the tilt stays under
    the hover's own maximum over the 5\,s before DETACH plus 0.5$^\circ$ (back in band); and the largest
    change of the payload's height from its mean over the 5\,s before, in the 12\,s after. On the rig, where
    the ring rocks in hover, the rise of the peak above its mean over the 3\,s before the release is also
    given.
    \item[Attach] The peak payload tilt and the height loss (dip) after the weld, and whether the tilt
    returns under 2$^\circ$ within 3\,s; for the rejoin mission, also the tilt at the detach, the mean tilt
    over the orbit and the time from TAKEOFF to landed.
    \item[Solver health] The number of planner solves that failed or did not converge.
\end{description}
Each table cell is the mean over its runs, with the number of runs in brackets.

% 11 Oct: the physical rationale for the bars (health check, item 3). The metric definitions themselves (the
% tilt-vector rise over a 5 s pre-window, the post-event band, event times) are frozen on Mon 12 per the thesis
% plan's T6 (P1, D4) and replace the detach item above once signed.
\paragraph{Why these bars.} Three physical limits say what a transient may do to the ring, and the bars of the
results sections sit inside them. The first is the magnets. A plate holds its magnet straight off the face by
about 9--10\,N and by less under shear, so a tilt that turns a rod's pull across its plate can let a magnet go:
every ring lost on the rig began with an uncommanded magnet release, all of them with the rods planned at
45$^\circ$, and the rods are planned at 60$^\circ$ since 7 October so that each pull stays closer to its plate's
axis (Section~\ref{sec:sim_detach}). [The hold against the angle of pull, measured with a spring scale at 0, 30,
45 and 60$^\circ$ off the axis on 14 October, gives this limit its number.] The second is the fleet's own rule:
a ring tilt over 30$^\circ$, or a second loss, drops the ring and lands every UAV (Section~\ref{sec:safety}), so
30$^\circ$ is where the controller gives up, not a value to design to. The third is the model: the planner's
load--rod model and each tracker's rod pull assume a near-level ring, so a transient that stays within a few
degrees stays inside the regime the plan was computed for, and a transient that leaves it is being flown on a
plan that no longer describes the ring. The bars are therefore stated as a rise above the ring's own motion
before the event, small against the first two limits and inside the third, and every run is reported with its
measured rise so that a run that misses a bar is still read against the limits rather than against a number.
